\chapter{Reproducibility}
\label{apx:reproducibility}

This appendix explains how to repeat the experiments.

\lstdefinestyle{shell}{language=bash, basicstyle=\footnotesize\ttfamily, frame=lines, columns=fullflexible,
  keepspaces=true, breaklines=true, numbers=none, xleftmargin=0pt, framexleftmargin=0pt,
  aboveskip=4pt, belowskip=4pt}

\section{Repository organization}
% Qué hay en cada carpeta del repositorio: implementación, guiones para el clúster, resultados y la tesis.

\begin{description}[nosep, itemsep=2pt]
  \item[\path{v2v/}] The implementation: the environment (\path{v2v_env/}); the agents, the two forms of
    aggregation, and the fixed rules (\path{agents/}); the training programs, \path{main.py} for the learning
    methods and \path{train_greedy.py} for the fixed rules; the analysis (\path{analysis.py},
    \path{plots.py}); and the tests (\path{tests/}).
  \item[\path{v2v/jobs/}] The scripts for the cluster (\path{*.sbatch}) and the options of each sweep
    (\path{sweeps/}) and of each variant of the diagnosis (\path{diag_variants/}).
  \item[\path{v2v/runs/}] One folder per sweep and value of $K$, with a \path{.csv} file per run and a
    \path{.json} file with its configuration, machine, and duration (the diagnosis, in \path{diag/}).
  \item[\path{Makefile}] The commands that submit the sweeps and bring the results back
    (\texttt{make pull-runs}).
  \item[\path{thesis-template/}] This document and the script that generates its tables and figures.
\end{description}

\section{Steps}
% Versiones, instalación y pruebas; una corrida; los barridos en el clúster; las tablas, las figuras y el documento.

\paragraph{Install and test.} The implementation uses Python~3.11 and the versions pinned in the lock file:
PyTorch~2.6.0, PettingZoo~1.27.0, Gymnasium~1.3.0, NumPy~2.4.6, pandas~3.0.5, Matplotlib~3.11.1, and
pytest~9.1.1. From the folder \path{v2v}, the package manager \texttt{uv} installs exactly those versions and
runs the tests:
\begin{lstlisting}[style=shell]
uv sync --frozen
uv run pytest
\end{lstlisting}

\paragraph{One run.} \looseness=-1 For example, the proposal with $K = 10$ and seed 0:
\begin{lstlisting}[style=shell]
uv run python main.py --aggregation gossip \
    --num-v2v-pairs 10 --episodes 2000 --seed 0 \
    --local-reward --reward-end-of-slot-queue \
    --gamma 0.9 --output runs/propia/k10/gossip_seed0.csv
\end{lstlisting}
The other methods use \texttt{--aggregation fedavg} or \texttt{none} (\acs{DTE}); the shared reward drops
\texttt{--local-reward} and \texttt{--reward-end-of-slot-queue}; the variant with $S = 2$ adds
\texttt{--gossip-neighbors 2}. The fixed rules run with \path{train_greedy.py}: Greedy by default, and the
threshold rule with \texttt{--policy umbral}.

\paragraph{The sweeps.} \looseness=-1 On the cluster, from the root of the repository, each sweep is submitted with one command
per method:
\begin{lstlisting}[style=shell]
make run-<method> SWEEP=propia RUN_ON=cpu
\end{lstlisting}
where the method is \texttt{fedavg}, \texttt{dte}, \texttt{gossip}, \texttt{greedy}, or \texttt{umbral} (the
last two without \texttt{RUN\_ON}). The sweeps keep their names in the repository: \texttt{propia} (individual
reward), \texttt{paper} (shared reward), and \texttt{propia-vecinos2} and \texttt{propia-vecinos3} (gossip with
$S = 2$ and $S = 3$). Each command submits one job per value of $K$, with seeds 0, 1, and 2, and skips the runs
that already have results. The diagnosis is submitted with \texttt{make run-diag DIAG\_SET=<set>}, where the set
is a file of \path{diag_variants/}.

\paragraph{The tables, the figures, and the document.} \looseness=-1 With the results in \path{v2v/runs/}, the script
\path{scripts/generate_results.py} of the thesis generates all the figures (\path{figures/res_*.pdf}) and tables
(\path{tables/res_*.tex}) of the results. It runs from the root of the repository with the Python environment of
the implementation (on Windows, \path{Scripts/python.exe} instead of \path{bin/python}):
\begin{lstlisting}[style=shell]
v2v/.venv/bin/python \
    thesis-template/scripts/generate_results.py
\end{lstlisting}
It stores what it reads in \path{scripts/_cache/}, which must be deleted if the runs change. Then
\texttt{latexmk -pdf main.tex}, in the folder of the thesis, compiles this document.

\section{Seeds}
% Cómo la semilla de una corrida determina posiciones, desvanecimientos, pesos iniciales y elección de vecinos, para que todos los métodos vean los mismos escenarios.

\looseness=-1 The seed $\varsigma$ of a run determines all its randomness, so two runs with the same seed
see the same scenarios, whatever the method (Section~\ref{sec:exp-runs}). Episode $e$ draws its positions
and its fading with the seed $10^6 \varsigma + e$, and agent $i$ draws its initial weights, its exploration, and
its samples from the buffer with the seed $10^6 \varsigma + i$. The choice of neighbors of the capped variant and
the tie-breaking of the fixed rules use generators of their own, derived from $\varsigma$. The exact results
can still differ between a processor and a graphics card, or between machines (Appendix~\ref{apx:diagnosis}).

\endinput
